% Ratio sweep: one control, f2/f1, crosses the line between harmonic and inharmonic.
% Every lattice line m + n·r (|m| + |n| <= 5) is straight in this view; at a simple ratio many meet and the
% spectrum folds onto one series. Above, the scan on the torus at three ratios: two close, the golden one never does.
% Data and geometry: ratio-sweep-gates.py (rsg-*.dat, rsg-torus-mesh.tex, rsg-knot-*.tex).
% Compile: lualatex ratio-sweep-gates.tex ; SVG via PyMuPDF (dvisvgm's PDF route is broken on this Mac).
\documentclass[border=0pt]{standalone}
\usepackage{fontspec}
\setmainfont{JetBrains Mono}
\setmonofont{JetBrains Mono}
\usepackage{tikz}
\usepackage{pgfplots}
\pgfplotsset{compat=1.16}
\usetikzlibrary{arrows.meta, calc, backgrounds}
% Drafting Table, paper ground (_ops/drafting-table/design-system/colors_and_type.css)
\definecolor{paper}{HTML}{F0EEE5}
\definecolor{ink}{HTML}{1A1916}
\definecolor{ink2}{HTML}{3A3A35}
\definecolor{ink3}{HTML}{6A6862}
\definecolor{slab}{HTML}{D8D4C6}
\tikzset{
  mesh/.style={draw=ink3!40, line width=0.2pt},
  knot/.style={draw=ink, line width=0.55pt},
  knotback/.style={draw=ink3!60, line width=0.3pt, densely dashed},
  lab/.style={font=\scriptsize, text=ink2},
  rowhead/.style={font=\bfseries\small, anchor=west, text=ink},
  rowsub/.style={font=\scriptsize, anchor=west, text=ink2},
  gate/.style={draw=ink2, line width=0.5pt, dash pattern=on 2.5pt off 1.5pt},
  lead/.style={draw=ink3, line width=0.3pt},
}
\pgfplotsset{/pgf/number format/assume math mode=true}
% x of a ratio on the page: the axis starts at 1.6 cm and is 17.4 cm wide for r in [1.30, 1.75]
\newcommand{\rx}[1]{\fpeval{1.6 + (#1 - 1.3) * 17.4 / 0.45}}
\begin{document}
\begin{tikzpicture}

% =============== the panel ===============
\draw[ink, line width=0.5pt] (-0.3,0) rectangle (19.9,-13.4);
\draw[ink, line width=0.3pt] (-0.3,-0.95) -- (19.9,-0.95);
\node[rowhead] at (0,-0.32) {SWEEP THE RATIO \enspace one knob, harmonic to inharmonic};
\node[rowsub] at (0,-0.7) {f₂ ÷ f₁ swept from 1.30 to 1.75; every line is a partial m·f₁ + n·f₂, with |m| + |n| ≤ 5};

% ---------- the strip: the scan on the torus at three ratios ----------
\foreach \nm/\x/\r/\cap in {32/9.33/1.5/{closes: a (2,3) knot}, phi/12.9/1.6180340/{φ: never closes}, 53/16.7/1.6666667/{closes: a (3,5) knot}} {
  \begin{scope}[shift={(\x,-1.72)}, scale=0.38]
    % mini torus mesh, generated by ratio-sweep-gates.py
\draw[mesh] (2.426,0.056) -- (2.415,0.111) -- (2.397,0.166) -- (2.371,0.218) -- (2.339,0.268) -- (2.300,0.315) -- (2.255,0.359) -- (2.205,0.398) -- (2.150,0.434) -- (2.090,0.464) -- (2.027,0.490) -- (1.960,0.510) -- (1.891,0.524) -- (1.821,0.533) -- (1.750,0.536) -- (1.679,0.533) -- (1.609,0.524) -- (1.540,0.510) -- (1.473,0.490) -- (1.410,0.464) -- (1.350,0.434) -- (1.295,0.398) -- (1.245,0.359) -- (1.200,0.315) -- (1.161,0.268) -- (1.129,0.218) -- (1.103,0.166) -- (1.085,0.111) -- (1.074,0.056) -- (1.070,0.000) -- (1.074,-0.056);
\draw[mesh] (1.992,1.023) -- (1.953,1.053) -- (1.910,1.077) -- (1.862,1.095) -- (1.810,1.107) -- (1.755,1.113) -- (1.698,1.113) -- (1.638,1.106) -- (1.577,1.093) -- (1.516,1.075) -- (1.454,1.050) -- (1.393,1.019) -- (1.334,0.984) -- (1.276,0.943) -- (1.221,0.898) -- (1.169,0.849) -- (1.121,0.797) -- (1.078,0.742) -- (1.039,0.684) -- (1.006,0.625) -- (0.978,0.565) -- (0.955,0.505) -- (0.940,0.445) -- (0.930,0.387) -- (0.927,0.329) -- (0.930,0.275) -- (0.940,0.223) -- (0.955,0.174) -- (0.978,0.130) -- (1.006,0.089) -- (1.039,0.054);
\draw[mesh] (1.103,1.574) -- (1.075,1.580) -- (1.045,1.578) -- (1.013,1.570) -- (0.980,1.555) -- (0.946,1.533) -- (0.911,1.504) -- (0.875,1.469) -- (0.839,1.428) -- (0.804,1.382) -- (0.770,1.331) -- (0.737,1.275) -- (0.705,1.216) -- (0.675,1.153) -- (0.647,1.089) -- (0.622,1.022) -- (0.600,0.955) -- (0.581,0.887) -- (0.564,0.820) -- (0.552,0.754) -- (0.542,0.690) -- (0.537,0.628) -- (0.535,0.571) -- (0.537,0.516) -- (0.542,0.467) -- (0.552,0.423) -- (0.564,0.384) -- (0.581,0.351) -- (0.600,0.325) -- (0.622,0.305) -- (0.647,0.292);
\draw[mesh] (0.000,1.757) -- (0.000,1.751) -- (0.000,1.737) -- (0.000,1.716) -- (0.000,1.689) -- (0.000,1.654) -- (0.000,1.613) -- (0.000,1.567) -- (0.000,1.515) -- (0.000,1.458) -- (0.000,1.397) -- (0.000,1.332) -- (0.000,1.265) -- (0.000,1.195) -- (0.000,1.125) -- (0.000,1.054) -- (0.000,0.983) -- (0.000,0.913) -- (0.000,0.845) -- (0.000,0.779) -- (0.000,0.717) -- (0.000,0.659) -- (0.000,0.605) -- (0.000,0.556) -- (0.000,0.514) -- (0.000,0.477) -- (0.000,0.447) -- (0.000,0.424) -- (0.000,0.408) -- (0.000,0.399) -- (0.000,0.398);
\draw[mesh] (-1.103,1.574) -- (-1.075,1.580) -- (-1.045,1.578) -- (-1.013,1.570) -- (-0.980,1.555) -- (-0.946,1.533) -- (-0.911,1.504) -- (-0.875,1.469) -- (-0.839,1.428) -- (-0.804,1.382) -- (-0.770,1.331) -- (-0.737,1.275) -- (-0.705,1.216) -- (-0.675,1.153) -- (-0.647,1.089) -- (-0.622,1.022) -- (-0.600,0.955) -- (-0.581,0.887) -- (-0.564,0.820) -- (-0.552,0.754) -- (-0.542,0.690) -- (-0.537,0.628) -- (-0.535,0.571) -- (-0.537,0.516) -- (-0.542,0.467) -- (-0.552,0.423) -- (-0.564,0.384) -- (-0.581,0.351) -- (-0.600,0.325) -- (-0.622,0.305) -- (-0.647,0.292);
\draw[mesh] (-1.992,1.023) -- (-1.953,1.053) -- (-1.910,1.077) -- (-1.862,1.095) -- (-1.810,1.107) -- (-1.755,1.113) -- (-1.698,1.113) -- (-1.638,1.106) -- (-1.577,1.093) -- (-1.516,1.075) -- (-1.454,1.050) -- (-1.393,1.019) -- (-1.334,0.984) -- (-1.276,0.943) -- (-1.221,0.898) -- (-1.169,0.849) -- (-1.121,0.797) -- (-1.078,0.742) -- (-1.039,0.684) -- (-1.006,0.625) -- (-0.978,0.565) -- (-0.955,0.505) -- (-0.940,0.445) -- (-0.930,0.387) -- (-0.927,0.329) -- (-0.930,0.275) -- (-0.940,0.223) -- (-0.955,0.174) -- (-0.978,0.130) -- (-1.006,0.089) -- (-1.039,0.054);
\draw[mesh] (-2.426,0.056) -- (-2.415,0.111) -- (-2.397,0.166) -- (-2.371,0.218) -- (-2.339,0.268) -- (-2.300,0.315) -- (-2.255,0.359) -- (-2.205,0.398) -- (-2.150,0.434) -- (-2.090,0.464) -- (-2.027,0.490) -- (-1.960,0.510) -- (-1.891,0.524) -- (-1.821,0.533) -- (-1.750,0.536) -- (-1.679,0.533) -- (-1.609,0.524) -- (-1.540,0.510) -- (-1.473,0.490) -- (-1.410,0.464) -- (-1.350,0.434) -- (-1.295,0.398) -- (-1.245,0.359) -- (-1.200,0.315) -- (-1.161,0.268) -- (-1.129,0.218) -- (-1.103,0.166) -- (-1.085,0.111) -- (-1.074,0.056) -- (-1.070,0.000) -- (-1.074,-0.056);
\draw[mesh] (-2.104,-0.748) -- (-2.101,-0.691) -- (-2.092,-0.632) -- (-2.076,-0.572) -- (-2.054,-0.512) -- (-2.026,-0.452) -- (-1.992,-0.393) -- (-1.953,-0.336) -- (-1.910,-0.281) -- (-1.862,-0.228) -- (-1.810,-0.179) -- (-1.755,-0.134) -- (-1.698,-0.094) -- (-1.638,-0.058) -- (-1.577,-0.028) -- (-1.516,-0.003) -- (-1.454,0.016) -- (-1.393,0.029) -- (-1.334,0.036) -- (-1.276,0.036) -- (-1.221,0.030) -- (-1.169,0.018) -- (-1.121,-0.000) -- (-1.078,-0.025) -- (-1.039,-0.054) -- (-1.006,-0.089);
\draw[mesh] (-2.026,-0.988) -- (-2.054,-0.948) -- (-2.076,-0.903) -- (-2.092,-0.855) -- (-2.101,-0.803) -- (-2.104,-0.748);
\draw[mesh] (-1.215,-1.296) -- (-1.213,-1.238) -- (-1.208,-1.176) -- (-1.198,-1.112) -- (-1.186,-1.046) -- (-1.169,-0.979) -- (-1.150,-0.911) -- (-1.128,-0.844) -- (-1.103,-0.777) -- (-1.075,-0.713) -- (-1.045,-0.650) -- (-1.013,-0.591) -- (-0.980,-0.535) -- (-0.946,-0.484) -- (-0.911,-0.438) -- (-0.875,-0.397) -- (-0.839,-0.362) -- (-0.804,-0.334) -- (-0.770,-0.311) -- (-0.737,-0.296) -- (-0.705,-0.288) -- (-0.675,-0.286) -- (-0.647,-0.292) -- (-0.622,-0.305);
\draw[mesh] (-1.128,-1.561) -- (-1.150,-1.541) -- (-1.169,-1.515) -- (-1.186,-1.482) -- (-1.198,-1.443) -- (-1.208,-1.399) -- (-1.213,-1.350) -- (-1.215,-1.296);
\draw[mesh] (-0.000,-1.496) -- (-0.000,-1.438) -- (-0.000,-1.375) -- (-0.000,-1.310) -- (-0.000,-1.242) -- (-0.000,-1.172) -- (-0.000,-1.101) -- (-0.000,-1.030) -- (-0.000,-0.959) -- (-0.000,-0.890) -- (-0.000,-0.823) -- (-0.000,-0.758) -- (-0.000,-0.697) -- (-0.000,-0.640) -- (-0.000,-0.588) -- (-0.000,-0.542) -- (-0.000,-0.501) -- (-0.000,-0.466) -- (-0.000,-0.438) -- (-0.000,-0.418) -- (-0.000,-0.404) -- (-0.000,-0.398) -- (-0.000,-0.399);
\draw[mesh] (-0.000,-1.756) -- (-0.000,-1.747) -- (-0.000,-1.731) -- (-0.000,-1.708) -- (-0.000,-1.678) -- (-0.000,-1.641) -- (-0.000,-1.598) -- (-0.000,-1.550) -- (-0.000,-1.496);
\draw[mesh] (1.215,-1.296) -- (1.213,-1.238) -- (1.208,-1.176) -- (1.198,-1.112) -- (1.186,-1.046) -- (1.169,-0.979) -- (1.150,-0.911) -- (1.128,-0.844) -- (1.103,-0.777) -- (1.075,-0.713) -- (1.045,-0.650) -- (1.013,-0.591) -- (0.980,-0.535) -- (0.946,-0.484) -- (0.911,-0.438) -- (0.875,-0.397) -- (0.839,-0.362) -- (0.804,-0.334) -- (0.770,-0.311) -- (0.737,-0.296) -- (0.705,-0.288) -- (0.675,-0.286) -- (0.647,-0.292) -- (0.622,-0.305);
\draw[mesh] (1.128,-1.561) -- (1.150,-1.541) -- (1.169,-1.515) -- (1.186,-1.482) -- (1.198,-1.443) -- (1.208,-1.399) -- (1.213,-1.350) -- (1.215,-1.296);
\draw[mesh] (2.104,-0.748) -- (2.101,-0.691) -- (2.092,-0.632) -- (2.076,-0.572) -- (2.054,-0.512) -- (2.026,-0.452) -- (1.992,-0.393) -- (1.953,-0.336) -- (1.910,-0.281) -- (1.862,-0.228) -- (1.810,-0.179) -- (1.755,-0.134) -- (1.698,-0.094) -- (1.638,-0.058) -- (1.577,-0.028) -- (1.516,-0.003) -- (1.454,0.016) -- (1.393,0.029) -- (1.334,0.036) -- (1.276,0.036) -- (1.221,0.030) -- (1.169,0.018) -- (1.121,-0.000) -- (1.078,-0.025) -- (1.039,-0.054) -- (1.006,-0.089);
\draw[mesh] (2.026,-0.988) -- (2.054,-0.948) -- (2.076,-0.903) -- (2.092,-0.855) -- (2.101,-0.803) -- (2.104,-0.748);
\draw[mesh] (-2.427,-0.078) -- (-2.417,-0.156) -- (-2.400,-0.234) -- (-2.377,-0.311) -- (-2.347,-0.387) -- (-2.311,-0.462) -- (-2.269,-0.536) -- (-2.220,-0.609) -- (-2.165,-0.679) -- (-2.104,-0.748) -- (-2.038,-0.815) -- (-1.966,-0.879) -- (-1.888,-0.941) -- (-1.806,-1.001) -- (-1.718,-1.058) -- (-1.626,-1.112) -- (-1.529,-1.163) -- (-1.428,-1.210) -- (-1.323,-1.255) -- (-1.215,-1.296) -- (-1.103,-1.333) -- (-0.988,-1.367) -- (-0.871,-1.397) -- (-0.751,-1.423) -- (-0.629,-1.445) -- (-0.505,-1.463) -- (-0.380,-1.478) -- (-0.254,-1.488) -- (-0.127,-1.494) -- (-0.000,-1.496) -- (0.127,-1.494) -- (0.254,-1.488) -- (0.380,-1.478) -- (0.505,-1.463) -- (0.629,-1.445) -- (0.751,-1.423) -- (0.871,-1.397) -- (0.988,-1.367) -- (1.103,-1.333) -- (1.215,-1.296) -- (1.323,-1.255) -- (1.428,-1.210) -- (1.529,-1.163) -- (1.626,-1.112) -- (1.718,-1.058) -- (1.806,-1.001) -- (1.888,-0.941) -- (1.966,-0.879) -- (2.038,-0.815) -- (2.104,-0.748) -- (2.165,-0.679) -- (2.220,-0.609) -- (2.269,-0.536) -- (2.311,-0.462) -- (2.347,-0.387) -- (2.377,-0.311) -- (2.400,-0.234) -- (2.417,-0.156) -- (2.427,-0.078) -- (2.430,-0.000);
\draw[mesh] (1.750,0.536) -- (1.748,0.592) -- (1.740,0.648) -- (1.728,0.704) -- (1.712,0.760) -- (1.690,0.815) -- (1.664,0.869) -- (1.634,0.922) -- (1.599,0.974) -- (1.559,1.025) -- (1.516,1.075) -- (1.468,1.123) -- (1.416,1.169) -- (1.360,1.214) -- (1.301,1.257) -- (1.237,1.298) -- (1.171,1.337) -- (1.101,1.373) -- (1.029,1.407) -- (0.953,1.439) -- (0.875,1.469) -- (0.794,1.496) -- (0.712,1.520) -- (0.627,1.542) -- (0.541,1.561) -- (0.453,1.577) -- (0.364,1.590) -- (0.274,1.600) -- (0.183,1.607) -- (0.092,1.612) -- (0.000,1.613) -- (-0.092,1.612) -- (-0.183,1.607) -- (-0.274,1.600) -- (-0.364,1.590) -- (-0.453,1.577) -- (-0.541,1.561) -- (-0.627,1.542) -- (-0.712,1.520) -- (-0.794,1.496) -- (-0.875,1.469) -- (-0.953,1.439) -- (-1.029,1.407) -- (-1.101,1.373) -- (-1.171,1.337) -- (-1.237,1.298) -- (-1.301,1.257) -- (-1.360,1.214) -- (-1.416,1.169) -- (-1.468,1.123) -- (-1.516,1.075) -- (-1.559,1.025) -- (-1.599,0.974) -- (-1.634,0.922) -- (-1.664,0.869) -- (-1.690,0.815) -- (-1.712,0.760) -- (-1.728,0.704) -- (-1.740,0.648) -- (-1.748,0.592) -- (-1.750,0.536) -- (-1.748,0.479) -- (-1.740,0.423) -- (-1.728,0.367) -- (-1.712,0.312) -- (-1.690,0.257) -- (-1.664,0.203) -- (-1.634,0.150) -- (-1.599,0.098) -- (-1.559,0.047) -- (-1.516,-0.003) -- (-1.468,-0.051) -- (-1.416,-0.097) -- (-1.360,-0.142) -- (-1.301,-0.185) -- (-1.237,-0.226) -- (-1.171,-0.265) -- (-1.101,-0.301) -- (-1.029,-0.336) -- (-0.953,-0.368) -- (-0.875,-0.397) -- (-0.794,-0.424) -- (-0.712,-0.448) -- (-0.627,-0.470) -- (-0.541,-0.489) -- (-0.453,-0.505) -- (-0.364,-0.518) -- (-0.274,-0.528) -- (-0.183,-0.536) -- (-0.092,-0.540) -- (-0.000,-0.542) -- (0.092,-0.540) -- (0.183,-0.536) -- (0.274,-0.528) -- (0.364,-0.518) -- (0.453,-0.505) -- (0.541,-0.489) -- (0.627,-0.470) -- (0.712,-0.448) -- (0.794,-0.424) -- (0.875,-0.397) -- (0.953,-0.368) -- (1.029,-0.336) -- (1.101,-0.301) -- (1.171,-0.265) -- (1.237,-0.226) -- (1.301,-0.185) -- (1.360,-0.142) -- (1.416,-0.097) -- (1.468,-0.051) -- (1.516,-0.003) -- (1.559,0.047) -- (1.599,0.098) -- (1.634,0.150) -- (1.664,0.203) -- (1.690,0.257) -- (1.712,0.312) -- (1.728,0.367) -- (1.740,0.423) -- (1.748,0.479) -- (1.750,0.536);
\draw[mesh] (1.070,0.000) -- (1.069,0.034) -- (1.064,0.069) -- (1.057,0.103) -- (1.047,0.137) -- (1.034,0.170) -- (1.018,0.204) -- (0.999,0.236) -- (0.977,0.268) -- (0.953,0.299) -- (0.927,0.329) -- (0.897,0.359) -- (0.866,0.387) -- (0.832,0.415) -- (0.795,0.441) -- (0.757,0.466) -- (0.716,0.490) -- (0.673,0.512) -- (0.629,0.533) -- (0.583,0.552) -- (0.535,0.571) -- (0.486,0.587) -- (0.435,0.602) -- (0.383,0.615) -- (0.331,0.627) -- (0.277,0.636) -- (0.222,0.644) -- (0.167,0.651) -- (0.112,0.655) -- (0.056,0.658) -- (0.000,0.659) -- (-0.056,0.658) -- (-0.112,0.655) -- (-0.167,0.651) -- (-0.222,0.644) -- (-0.277,0.636) -- (-0.331,0.627) -- (-0.383,0.615) -- (-0.435,0.602) -- (-0.486,0.587) -- (-0.535,0.571) -- (-0.583,0.552) -- (-0.629,0.533) -- (-0.673,0.512) -- (-0.716,0.490) -- (-0.757,0.466) -- (-0.795,0.441) -- (-0.832,0.415) -- (-0.866,0.387) -- (-0.897,0.359) -- (-0.927,0.329) -- (-0.953,0.299) -- (-0.977,0.268) -- (-0.999,0.236) -- (-1.018,0.204) -- (-1.034,0.170) -- (-1.047,0.137) -- (-1.057,0.103) -- (-1.064,0.069) -- (-1.069,0.034) -- (-1.070,0.000) -- (-1.069,-0.034);

    \input{rsg-knot-\nm.tex}
  \end{scope}
  \node[lab, anchor=north] (c\nm) at (\x,-2.42) {\cap};
  \draw[lead] (c\nm.south) -- (\rx{\r},-3.12);
}
% ---------- gate names, just above the plot ----------
\foreach \r/\t in {1.3333333/4\slash 3, 1.5/3\slash 2, 1.6/8\slash 5, 1.6180340/φ, 1.6666667/5\slash 3}
  \node[font=\scriptsize\bfseries, text=ink, anchor=south] at (\rx{\r},-3.42) {\t};

\begin{axis}[at={(1.6cm,-12.0cm)}, scale only axis, width=17.4cm, height=8.4cm,
    xmin=1.3, xmax=1.75, ymin=0, ymax=6, axis lines=left, tick align=outside, clip=false,
    xtick={1.3,1.4,1.5,1.6,1.7}, xticklabels={1.30,1.40,1.50,1.60,1.70}, ytick={0,1,2,3,4,5,6},
    axis line style={ink, line width=0.5pt}, tick style={ink},
    xlabel={ratio r = f₂ ÷ f₁}, ylabel={frequency ÷ f₁},
    xlabel style={font=\scriptsize, text=ink2, yshift=1mm}, ylabel style={font=\scriptsize, text=ink2},
    tick label style={font=\tiny, text=ink2}, empty line=jump]
  % gates, under everything
  \foreach \r in {1.3333333, 1.5, 1.6, 1.6180340, 1.6666667} { \edef\t{\noexpand\draw[gate] (axis cs:\r,0) -- (axis cs:\r,6.06);} \t }
  % the lattice, fading with order |m| + |n|
  \addplot[ink, opacity=0.22, line width=0.35pt] table[x=r, y=f] {rsg-lines-w5.dat};
  \addplot[ink, opacity=0.35, line width=0.4pt] table[x=r, y=f] {rsg-lines-w4.dat};
  \addplot[ink, opacity=0.55, line width=0.45pt] table[x=r, y=f] {rsg-lines-w3.dat};
  \addplot[ink, opacity=0.75, line width=0.55pt] table[x=r, y=f] {rsg-lines-w2.dat};
  \addplot[ink, line width=0.75pt] table[x=r, y=f] {rsg-lines-w1.dat};
  % folds: where two or more drawn lines meet on a gate
  \addplot[only marks, mark=*, mark size=1.5pt, mark options={draw=ink, fill=paper, line width=0.5pt},
    restrict expr to domain={\thisrow{count}}{2:2}] table[x=r, y=f] {rsg-folds.dat};
  \addplot[only marks, mark=*, mark size=2.1pt, mark options={draw=ink, fill=ink},
    restrict expr to domain={\thisrow{count}}{3:3}] table[x=r, y=f] {rsg-folds.dat};
  \node[lab, anchor=west] at (axis cs:1.75,1.75) {f₂};
  \node[lab, anchor=west] at (axis cs:1.75,1) {f₁};
\end{axis}
\node[lab, anchor=north east, align=right] at (19.7,-12.45) {○ two lines meet \quad ● three lines meet};

\node[font=\scriptsize, anchor=north west, text=ink2, text width=19.6cm] at (0,-13.65) {\textbf{At a simple ratio the scan closes and the lattice folds onto one series; just off it, each partial slides away slowly.} At 3/2 every drawn line lands on a multiple of f₁/2. At 8/5 no two drawn lines meet: that gate folds only at higher orders. At φ none ever do; the nearest two drawn lines stay 0.146 f₁ apart.};

% ---------- title block ----------
\draw[ink, line width=0.8pt] (-0.3,-14.9) rectangle (19.9,-15.8);
\draw[ink, line width=0.3pt] (9.6,-14.9) -- (9.6,-15.8) (15.3,-14.9) -- (15.3,-15.8);
\node[font=\scriptsize\bfseries, text=ink, anchor=west] at (-0.1,-15.35) {DSP IN LOOPING DIMENSIONS · RATIO SWEEP GATES};
\node[font=\scriptsize, text=ink, anchor=west] at (9.8,-15.35) {INK computed};
\node[font=\scriptsize, text=ink2, anchor=west] at (15.5,-15.35) {FIG · 2026-11-09 · PAPER};

\begin{scope}[on background layer]
  \fill[paper] ($(current bounding box.south west)+(-0.4,-0.4)$) rectangle ($(current bounding box.north east)+(0.4,0.4)$);
\end{scope}
\end{tikzpicture}
\end{document}
